% REVISION NOTES (highlighted content marked with %[NEW] or %[CHANGED])
% Changes made in response to reviewer comments:
%   - Reviewer B: added HMM+causal discovery (Balsells2024), distribution
%     shift context (Mameche2025), recent neural Granger (Marcinkevics2021,
%     Qian2026), J-PCMCI+ (Guenther2023)
%   - Reviewer A: added LLMs in scientific ML subsection (Wuwu2025, Tiwari2025)
%   - Editorial: replaced 6 pre-2021 refs with 8 refs from 2021-2026 to
%     achieve >= 80% recency requirement
%   - Critical LLM perspective added (Long2023, Wu2025)
% All new or substantially changed sentences are marked with \rev{...}
% Requires: \highlightchangestrue/false toggle defined in preamble

\section{Related Work}
\label{sec:related}

\subsection{Classical Causal Discovery and Continuous Structure Learning}

Score-based causal discovery methods recover directed acyclic graphs
by optimising a statistical criterion over the space of DAG structures.
A major algorithmic breakthrough came with NOTEARS~\cite{Zheng2018},
which reformulated structure learning as a fully continuous optimisation
problem by introducing the smooth acyclicity constraint
$h(\mathbf{W}) = \mathrm{tr}(\exp(\mathbf{W} \circ \mathbf{W})) - d = 0$,
enabling gradient-based DAG learning over the full adjacency matrix and
making downstream integration with neural architectures tractable.
REIGN inherits this acyclicity regulariser and applies it jointly with
the neural Granger and LLM-prior objectives inside each regime window.
\rev{Despite its elegance, the NOTEARS framework assumes stationarity of the
underlying data-generating process, a restriction that is rarely satisfied
in real-world financial time series exhibiting regime-switching
dynamics~\cite{Mameche2025,Balsells2024}.}                  %[NEW sentence]

\subsection{Neural Granger Causality}

Classical Granger causality tests each variable pair in isolation,
making it susceptible to spurious edges from indirect effects and shared
latent drivers.
Tank et al.~\cite{Tank2022} introduced component-wise MLPs and LSTMs
(cMLP, cLSTM) whose group-lasso-penalised first-layer input weights
serve as Granger causality indicators, establishing the paradigm of
end-to-end neural Granger learning from multivariate time series.
\rev{Marcinkev\v{c}s and Vogt~\cite{Marcinkevics2021} extended this paradigm
with Generalised VAR (GVAR), a self-explaining neural network that jointly
infers Granger-causal structure and detects the sign and time-varying
magnitude of each causal effect, yielding substantially more interpretable
outputs than sparse-input networks.}                         %[NEW sentence]
CUTS~\cite{Cheng2023} extended neural Granger causality to irregular time
series through an alternating imputation-and-discovery scheme using a
delayed-supervision GNN, directly addressing the irregular-sampling
challenge common in business KPI datasets.
CUTS+~\cite{Cheng2024} further scaled this framework to
high-dimensional settings via a coarse-to-fine discovery (C2FD) procedure
and a message-passing GNN, and constitutes the primary backbone of REIGN.
\rev{Most recently, Qian et al.~\cite{Qian2026} proposed MSDG, a multiscale
dynamic graph neural network that reconstructs multivariate time series
into dynamic graphs using causal windows of varying lengths and applies a
graph attention network jointly with an LSTM to capture time-varying
causal contributions; while MSDG advances dynamic Granger learning,
it relies on heuristic window selection and does not leverage LLM semantic
priors or statistically principled regime detection.}        %[NEW sentence]
Our framework extends CUTS+ along three orthogonal dimensions:
LLM-prior regularisation, regime-aware segmentation, and
confidence-weighted ensemble fusion, none of which appear in any of
the above designs.

\subsection{Time-Series Causal Discovery}

Time-series causal discovery must additionally account for temporal
autocorrelation, lagged dependencies, and the possibility of
contemporaneous effects.
Runge~\cite{Runge2020} proposed PCMCI+, which separates the lagged and
contemporaneous edge-removal phases and uses momentary conditional
independence to control for autocorrelation, achieving substantially
higher recall than constraint-based predecessors on observational time
series.
\rev{G\"unther et al.~\cite{Guenther2023} proposed J-PCMCI+, which
extends PCMCI+ to jointly infer causal graphs across multiple
heterogeneous time-series datasets that share latent contextual variables,
enabling discovery under cross-dataset distribution shift; this setting
closely resembles real-world business deployments where the same
causal graph applies across but differs between organisational units or
product lines.}                                              %[NEW sentence]
\rev{Mameche et al.~\cite{Mameche2025} introduced SPACETIME, a functional
non-parametric framework that jointly reconstructs temporal regimes and
causal mechanisms using Gaussian processes and minimum-description-length
scoring, explicitly unifying the tasks of regime detection, dataset
partitioning, and causal graph discovery under nonstationarity; SPACETIME
demonstrates that treating nonstationarity as an opportunity rather than
a nuisance consistently improves edge precision on climate and ecological
datasets, a finding we exploit in REIGN's regime-local design.}
                                                             %[NEW sentence]

\subsection{LLM-Augmented Causal Discovery}

Integrating large language models into causal discovery pipelines has
emerged as a rapidly developing research direction, motivated by the
observation that LLMs encode rich, text-derived domain knowledge about
cause-effect relationships.
Jin et al.~\cite{Jin2024ICLR} established the foundational result
that LLMs can serve as effective probabilistic priors for causal graphs
by querying pairwise edge existence and direction to produce soft
adjacency probabilities compatible with score-based discovery pipelines.
\rev{Long et al.~\cite{Long2023} investigated how LLMs can be treated as
imperfect domain experts whose erroneous causal claims must be corrected
by enforcing acyclicity and Markov-equivalence-class consistency before
integration into discovery pipelines; their analysis of failure modes
directly motivates REIGN's use of confidence-weighted soft priors rather
than hard constraints.}                                      %[NEW sentence]
Du et al.~\cite{Du2025} proposed LLM-CD, which integrates LLM-derived
knowledge at multiple stages of the PC algorithm, reporting an average
recall improvement of 169\% and a maximum of 403\% on the WUSCH dataset
compared to purely data-driven baselines.
Roy et al.~\cite{Roy2025} introduced Causal-LLM, a unified one-shot
framework combining prompt-based and data-driven modes that outperforms
classical score-based methods by approximately 40\% in edge accuracy
on the Asia and Sachs benchmarks.
Kampani et al.~\cite{Waxman2024} proposed LLM-initialised differentiable
causal discovery (LLM-DCD), using LLM-generated soft adjacency as a
warm-start initialisation for gradient-based DAG learning, the published
architecture most similar to REIGN; our framework differs by employing
neural Granger scoring throughout training and by applying regime
segmentation prior to discovery.
\rev{Liu et al.~\cite{Liu2026} introduced LLM-GC, the first method to
directly bridge LLMs and neural Granger causality for time series via
dual-modality encoding and a causality-aware self-attention mechanism;
however, LLM-GC does not address nonstationarity, regime switching, or
DAG acyclicity enforcement, all of which are central to REIGN.}
                                                             %[NEW sentence]
Vashishtha et al.~\cite{Vashishtha2025} demonstrated that eliciting a
topological causal ordering from an LLM, rather than pairwise edge scores,
yields more reliable priors by suppressing cyclic responses inherent in
independent pair-wise prompting; REIGN adopts this causal-order correction
step to post-process $\mathbf{A}_{\mathrm{LLM}}$ before fusing it into
the CUTS+ training objective.
\rev{Ze\v{c}evi\'{c} et al.~\cite{Zecevic2023} critically re-evaluated the role of LLMs in
causal discovery, arguing that their autoregressive, correlation-driven
modelling inherently lacks the theoretical grounding for causal reasoning,
and demonstrated through empirical studies that deliberate prompt
engineering can overstate LLM-based causal discovery performance by
injecting ground-truth knowledge; this finding reinforces REIGN's design
choice of restricting LLMs to a non-decisional soft-prior role whose
influence is continuously modulated by data evidence through the
regularisation parameter $\lambda$.}                        %[NEW sentence]

\subsection{Nonstationary and Regime-Aware Causal Discovery}

Financial and business time series exhibit well-documented
nonstationarity: the causal graph governing customer churn during a
promotional campaign window differs structurally from that in a
steady-state period.
\rev{Balsells-Rodas et al.~\cite{Balsells2024} established identifiability
guarantees for regime-dependent causal discovery using high-order Markov
Switching Models (MSMs), which extend Hidden Markov Models by
incorporating autoregressive dependencies on observations given discrete
latent states; their theoretical framework provides the foundational
justification for REIGN's assumption that locally stationary causal
structure can be recovered within each BOCPD/PELT-detected segment.}
                                                             %[NEW sentence]
\rev{Mameche et al.~\cite{Mameche2025} complemented this identifiability
analysis with an empirical demonstration of joint regime-and-graph
recovery in a non-parametric functional framework, confirming that
regime-aware methods consistently outperform single-graph alternatives
under real-world temporal distribution shift.}               %[NEW sentence]
To identify regime boundaries, REIGN supports both the PELT
algorithm~\cite{Killick2012} for offline retrospective discovery and
Bayesian Online Changepoint Detection (BOCPD)~\cite{Deasy2023} for
prospective deployment.
PELT provides exact dynamic-programming detection with linear expected
runtime, while BOCPD yields a principled probabilistic run-length
posterior that allows REIGN to quantify boundary uncertainty and propagate
it into the confidence-weighted ensemble.

\subsection{LLMs in Scientific and Data-Driven Machine Learning} %[NEW subsection]

\rev{The success of LLMs in causal discovery reflects a broader trend of
leveraging large foundation models to guide data-driven scientific
modelling, providing a methodological context for REIGN's prior-injection
design.
Wuwu et al.~\cite{Wuwu2025} introduced PINNsAgent, an LLM-based
agent that automates physics-informed neural network architecture
selection and hyperparameter tuning for partial differential equations
(PDEs) by encoding prior knowledge of solved PDEs as structured memory
and performing Memory Tree Reasoning over the design space; PINNsAgent
exemplifies how domain knowledge encoded in LLMs can steer a
computationally expensive search process without supplanting physical
modelling, a principle that directly parallels REIGN's use of LLM-derived
priors to regularise, rather than replace, neural Granger estimation.
Tiwari et al.~\cite{Tiwari2025} proposed the Latent Mamba Operator
(LaMO), which integrates state-space model efficiency in latent space
with the expressive power of kernel integral formulations to solve PDEs
on diverse geometries; LaMO's architecture, combining a compact latent
representation with a structured long-range dependency mechanism, shares
the computational philosophy of CUTS+'s coarse-to-fine discovery, where
a cheap first pass reduces the effective dimensionality before fine-grained
estimation.}                                                %[NEW paragraph]

\subsection{Evaluation Benchmarks}

Sachs et al.~\cite{Sachs2005} produced the canonical real-world
benchmark for causal discovery: an 11-node protein-signalling network
derived from 853 single-cell flow-cytometry measurements under known
interventions, against which the static slice of REIGN's output is
validated.
Stein et al.~\cite{Stein2025} released CausalRivers, the largest
in-the-wild time-series causal discovery benchmark to date, comprising
over 1,000 river-discharge stations at 15-minute resolution with curated
ground-truth causal graphs; the authors find that Granger-based approaches
remain competitive on this real-world benchmark, directly supporting the
choice of neural Granger as REIGN's backbone.
Ferdous et al.~\cite{Ferdous2025} introduced TimeGraph, a suite of
synthetic time-series datasets with independently controllable violation
factors including nonstationarity, irregular sampling, missing data, and
heterogeneous noise, enabling fine-grained ablation studies of each
REIGN component under isolated assumption violations.
Together, these three benchmarks form the multi-tier evaluation protocol
adopted in this paper.

% ============================================================
% REIGN Section 2.8 — Recent Developments (2022-2026)
% New subsection to be inserted after Section 2.7
% All new  keys defined in section28_refs.bib
% Papers marked [VERIFY] need author confirmation before submission
% ============================================================

% Section 2.8 — Insert after Section 2.7
% All \cite{} keys defined in section28_refs.bib
% [VERIFY] entries need author confirmation before submission

\subsection{Recent Developments in Causal Discovery and
             Foundation Models (2022--2026)}
\label{sec:recent}

\noindent\textbf{Neural causal discovery.}
Gong et al.\ \cite{Gong2023Rhino} introduced Rhino, combining
vector auto-regression, deep learning, and variational inference
to model history-dependent noise in temporal causal discovery.
Lippe et al.\ \cite{Lippe2022CITRIS} established identifiability
conditions for causal representations from temporal intervention
sequences, and Annadani et al.\ \cite{Annadani2023BayesDAG}
proposed BayesDAG for full Bayesian posterior inference over DAGs.
Geffner et al.\ \cite{Geffner2022DECI} developed DECI for
joint causal discovery and treatment effect estimation, while
Ashman et al.\ \cite{Ashman2023ADMG} addressed latent confounders
via neural ADMG learning.
Assaad et al.\ \cite{Assaad2022Survey} surveyed and evaluated
causal discovery methods for both i.i.d.\ and time-series data,
confirming that PCMCI+ and neural Granger methods are most
reliable under autocorrelation.
Cheng et al.\ \cite{Cheng2024CausalTime} introduced CausalTime,
a pipeline for generating realistic time series with ground-truth
causal graphs from real observational data.
Further advances include amortised variational
inference~\cite{Lorch2022AVICI}, diffusion-based structure
learning~\cite{Sanchez2023DiffAN}, and gradient-based nonlinear
Granger estimation~\cite{Tank2022}.

\noindent\textbf{LLMs and causal reasoning.}
Kiciman et al.\ \cite{Kiciman2024TMLR} benchmarked LLM causal
reasoning, finding GPT-4 achieves 97\% on pairwise causal
discovery and 92\% on counterfactual tasks, justifying the use
of LLM-derived structural priors in REIGN.
Zhang et al.\ \cite{Zhang2024Survey} and Jin et al.\
\cite{Jin2024ICLR} surveyed and empirically evaluated LLM-causal
integration, identifying prior injection as the most productive
mode while flagging prompt sensitivity and hallucination under
distributional shift as key risks.
A complementary IJCAI 2025 survey~\cite{Survey2025IJCAI}
catalogues regime-aware temporal causal discovery as an
underexplored gap that REIGN directly addresses.
Recent work has further studied LLMs as autonomous causal
agents~\cite{Jin2024EconAgent}, LLM understanding of causal
graph structure~\cite{Chen2024CLEAR}, interventional reasoning
evaluation~\cite{Kasetty2024Interventional}, and causal
reasoning fallacies~\cite{Joshi2024Fallacies}.

\noindent\textbf{Neural operators and scientific foundation models.}
The LLM-guided scientific ML community has developed
complementary architectures that share REIGN's principle of
structured knowledge guiding data-driven estimation.
Cheng et al.\ \cite{Cheng2026MNO} formally connected Mamba
state-space models to neural operators (MNO), demonstrating
superior long-range dependency capture over Transformers.
Jiang~\cite{Jiang2026AFDONet} and the AMO team~\cite{AMO2026ICLR}
applied adaptive Fourier decomposition theory to neural operator
design for inverse PDE problems and arbitrary-geometry meshes.
Additional advances include geometry-aware Mamba
operators~\cite{Gu2024Mamba}, physics-informed operator
learning~\cite{Zanardi2023PhysicsInformed}, and neural operators
for PDE learning~\cite{Kovachki2023}.

\noindent\textbf{Temporal nonstationarity and financial applications.}
The toolkit for nonstationary causal discovery has expanded with
methods for periodic nonstationarity~\cite{Gao2024PCMCI},
proxy-variable conditional independence testing~\cite{Miao2024Proxy},
self-compatibility evaluation without ground truth~\cite{Mogensen2024Compat},
nonstationary sparse transition representation learning~\cite{Chen2024CtrlNS},
and causal foundation models for time series~\cite{Wen2023Foundation}.
Score-based approaches using neural ODEs~\cite{Bellot2022NeurODE},
single-variable interventions for causal ordering~\cite{Chevalley2024CD},
and Transformer-based temporal discovery~\cite{Chen2023Transformer}
extend the differentiable structure learning paradigm.
In financial applications, regime-dependent Granger networks
have been shown to outperform single-graph alternatives for
systemic risk modelling~\cite{Billio2024Spillover}, causal churn
modelling advances interventional
analysis~\cite{Ouyang2024Churn}, and the MSM identifiability
framework has been extended to instantaneous effects and
exponential family noise~\cite{Dong2023MSM}.